\subsubsection{Geração de código}

A geração de código é tarefa de 224 estudos. As sínteses parciais cobrem 222 deles, divididos pela técnica focal em três grupos que não se sobrepõem:
\begin{itemize}
\item 92 estudos cuja técnica acrescenta \textit{feedback}, arranjo multiagente ou amostragem de candidatos;
\item 52 com raciocínio ou decomposição, sem esses mecanismos;
\item 78 apenas com instrução, exemplos, contexto, restrições ou otimização do prompt.
\end{itemize}

Os agrupamentos a seguir atravessam os grupos, e um estudo pode aparecer em mais de um. As contagens somam as sínteses parciais de cada grupo; como cada síntese delimitou o agrupamento à sua maneira, as somas são aproximadas e vêm marcadas como ``nas sínteses parciais''. As fichas foram extraídas por IA e ainda não foram conferidas por humano; a qualidade também foi julgada por IA. As comparações foram escolhidas por regra (resultado principal declarado, primeira tabela ou cobertura), e muitas técnicas são propostas dos próprios autores; a proporção de melhoras, portanto, não estima a eficácia geral de uma técnica.

\paragraph{Autorreparo e \textit{feedback} de execução}
Avaliam-se laços em que mensagens do compilador, testes falhos, saídas do interpretador ou alertas de analisadores voltam ao prompt para regenerar o código. Nas sínteses parciais, cerca de 37 estudos entram aqui. A maior parte dos contrastes é de pacote; o laço é isolado como ingrediente numa minoria (entre outros, S150, S479, S645, S878, S1047, S1086, S1898 e S1977), e S917, S1833, S2055 e S2163 são cumulativos.

% AFIRMACAO: A61
Contra um prompt de chamada única, os pacotes de reparo elevaram o Pass@1 de 73,78\% para 90,24\% no HumanEval com gpt-3.5-turbo-0613 ($\mu$FiX; S150-C1) e de 8,73\% para 28,0\% no CoderEval-Python (CoCoGen; S302-C1), e a retirada da fase de \textit{feedback} do $\mu$FiX deixou o valor em 85,37\% (S150-C2).
% AFIRMACAO: A62
Em cinco comparações de ingrediente com valores absolutos, o resultado com o laço foi maior, em métricas que não se agregam: erros de execução de 9 para 6 (S878-C1), HMean \textit{accuracy} de 12,9\% para 18,3\% (S1047-C3), acurácia de 61,7\% para 65,0\% (S1053-C3), Pass@1 de 0,90 para 0,94 sobre modelo ajustado (S1077-C1) e acurácia de 19,6\% para 22,6\% (S1086-C6).
% AFIRMACAO: A63
Nos dois estudos que igualam o orçamento de tentativas ou de programas, não houve ganho: no CodeContests, com pass 1@3, as variantes multiturno com \textit{feedback} ficaram em 2,4\% e 2,8\% contra 2,9\% com Llama 3.0 8B e em 6,9\% contra 7,7\% com Llama 3.1 8B (S479-C2 a S479-C4), e o SEIDR resolveu em média 11,5 problemas do PSB2, contra 12,2 da amostragem sem reparo e 12,0 da busca só de reparo (S645-C1, S645-C2).
% AFIRMACAO: A64
Com Llama 3.2 3B, pôr os testes no prompt rendeu acurácia maior (48,2\%) que o laço de interpretador isolado (22,6\%) e que a combinação dos dois (42,2\%), contra 19,6\% da base (S1086-C5, S1086-C6, S1086-C1).
% AFIRMACAO: A65
Num dos poucos contrastes que isolam a fonte do \textit{feedback}, com 10 requisitos, o relatório estático estruturado e a autorrevisão genérica produziram, cada um, 27 saídas válidas em 30 (S1763-C3).

O laço depende de testes, executor ou compilador disponíveis. Fora de S479 e S645, nenhuma comparação iguala o número de chamadas; uma leitura possível é que parte dos ganhos de pacote acompanhe o acréscimo de chamadas, o que os estudos não isolam. No esclarecimento de requisitos, o usuário simulado vê os testes de referência, condição oráculo (S2140-C2).

\paragraph{Autocrítica e revisão sem execução}
Avalia-se a crítica da saída pelo próprio modelo, ou por relatório externo que não executa o código, seguida de reescrita, muitas vezes dirigida a viés, segurança ou tratamento de exceções. Nas sínteses parciais, são 9 estudos (S114, S140, S149, S397, S604, S759, S834, S842 e S1204), quase todos com contraste de pacote.

% AFIRMACAO: A66
Com PaLM-2-CodeChat-bison, as cinco estratégias de mitigação de viés sem \textit{feedback} elevaram o CBS (menor é melhor) de 17,96\% para valores entre 23,05\% e 33,83\%, e as cinco com relatório de testes de viés o reduziram para valores entre 3,89\% e 6,29\% (S140-C1 a S140-C10).
% AFIRMACAO: A67
A autocrítica RCI reduziu a densidade de fraquezas apontadas pelo Bandit de 0,103 para 0,029 por linha de código com GPT-3 text-davinci-002 (S397-C1) e a proporção de amostras vulneráveis de 5,74\% para 4,5\% com GPT-3.5-turbo (S604-C11), mas, em S397, o teste de Kruskal-Wallis sobre a densidade não foi significativo para GPT-3 (p = 0,334) nem para GPT-3.5 (p = 0,160) (S397-C1).
% AFIRMACAO: A68
Com o modelo pequeno Qwen2.5-1.5B, Self-Refine (38,2\%) e SETS (38,8\%) tiveram Accuracy@t2 menor que a correção simples (44,6\%) (S759-C1, S759-C3), enquanto, em modelos maiores, os ganhos registrados são pequenos, como o CodeBLEU de 0,732 para 0,743 com Gemini-2.0 (S842-C5) e a completude de 92,9\% para 96,0\% com GPT-5 (S1204-C2).

S140 usa um único modelo, e seus autores admitem que parte da redução pode vir só da mudança do prompt. A dependência do porte é hipótese desta revisão, não teste. Analisadores estáticos e juízes LLM não demonstram segurança em uso real.

\paragraph{Sistemas multiagente}
Avaliam-se encadeamentos de agentes com papéis (analista, planejador, testador, revisor) e arranjos em que um modelo externo critica ou implementa. Nas sínteses parciais, são cerca de 29 estudos, todos avaliados como pacote, com ablações de componentes em cerca de dez. Nenhuma comparação com valores contrasta multiagente e agente único com o mesmo número de chamadas.

% AFIRMACAO: A69
Os pacotes elevaram o Pass@1 de 48,1\% para 80,5\% (MapCoder; S344-C1), de 67,68\% para 87,8\% (PairCoder; S448-C1) e de 57,3\% para 74,4\% (Self-collaboration; S2131-C1), e, em S1747, o MapCoder passou de 39,63\% para 64,02\% com Llama3.1 8B usando em média 8,1 chamadas contra 1,0 da base (S1747-C3).
% AFIRMACAO: A70
Em S344, S448, S530 e S634, as oito ablações que retiram um componente registraram valor menor que o do sistema completo, cada uma num único modelo e \textit{benchmark} (S344-C2, S344-C3, S448-C2, S448-C3, S530-C2, S530-C3, S634-C2, S634-C3).
% AFIRMACAO: A71
Há também resultados contrários e nulos: com CodeLlama 34B, o Agents4PLC compilou 31,3\% contra 68,8\% do LLM4PLC (S492-C1); substituir o módulo de codificação progressiva do Spec2RTL-Agent manteve 3 de 3 documentos corretos (S795-C3); e, com Qwen2.5-1.5B, o External-Refine guiado por GPT-4.1 ficou em 41,8\% contra 44,6\% (S759-C2).

Como prompting e orquestração variam juntos, o ganho não se atribui ao mecanismo multiagente. Com modelo externo mais forte (S759-C4, S1204-C3), o efeito se confunde com a seleção de modelo. Há ajuste fino de agentes em S1815, S1828 e S2170, e a única ablação do arranjo multiagente (S530-C3) usa \textit{benchmark} próprio.

\paragraph{Amostragem e seleção de candidatos}
Avalia-se a geração de várias candidatas com escolha por testes, ranqueamento, modelo de valor ou busca em árvore. Nas sínteses parciais, são cerca de 17 estudos, com ingrediente isolado em parte deles (S301, S399, S416, S448, S582, S645, S724 e S1034).

% AFIRMACAO: A72
Em pacotes, o VRank elevou o Pass@1 de 48,8\% para 58,1\% no VerilogEval-Human com GPT-4o-mini (S582-C1) e o SPECFIX, de 87,8\% para 92,0\% no HumanEval+ com DeepSeek-V3 (S724-C1); na ablação deste, retirar a entropia semântica ou a consistência com exemplos levou o valor a 90,55\% e 89,95\%, contra 91,99\% do sistema completo (S724-C2, S724-C3).
% AFIRMACAO: A73
Entre os contrários e nulos, a amostragem com $k = 5$ baixou a criatividade convergente de 2,52 para 2,02 com GPT-3.5 (S399-C3), e a retirada do CoT adaptativo do SEER mudou o pass@1 de 82,0\% para 81,0\%, diferença menor que as dispersões impressas (S1034-C3).

Em S399, a seleção usa o próprio oráculo de avaliação, indisponível sem testes de referência. O SEER inclui modelo de valor treinado, e o TIPCODER, proponente ajustado por SFT e RL; nesses casos (S1034, S1592), o contraste não mede só prompting.

\paragraph{Instrução de raciocínio}
Avalia-se a instrução para raciocinar passo a passo antes do código, ou a inserção de uma cadeia de raciocínio, sobre um prompt que, fora isso, fica igual. Nas sínteses parciais, são cerca de 14 estudos, a maioria com contraste de ingrediente.

% AFIRMACAO: A74
No grupo com raciocínio ou decomposição, as 8 comparações, em 8 estudos, que isolam a instrução de raciocínio como ingrediente dividem-se em 4 melhoras e 4 pioras (S124-C2, S207-C1, S309-C2, S319-C1, S491-C3, S508-C1, S560-C2, S569-C1), e S560 relata ausência de efeito significativo do CoT na correção nos três modelos avaliados (S560-C2).
% AFIRMACAO: A75
No grupo com \textit{feedback}, multiagente ou amostragem, o CoT reduziu o pass@1 de 18,24\% para 14,86\% com Qwen2.5-Coder-7B no BigCodeBench (S1121-C1), de 70,87\% para 61,11\% sobre o prompt EAS com DeepSeek-R1-Distill-Qwen-14B (S1024-C4) e a correção média de 0,53 para 0,49 com Llama-3.1-8B, sob um limite de 512 \textit{tokens} gerados que o estudo associa ao truncamento (S1606-C4), e a elevou de 81,0\% para 87,68\% com DeepSeek-V3 num \textit{benchmark} customizado (S1229-C2).
% AFIRMACAO: A76
Quando a cadeia veio de outro modelo, o pass@1 do Qwen2.5-Coder-7B no BigCodeBench subiu de 18,24\% para 28,38\% com o raciocínio do DeepSeek-R1 (S1121-C4), contraste que não separa o raciocínio do modelo que o produziu.

Comparações de uma célula (um modelo e um \textit{benchmark}) em \textit{benchmarks} distintos não separam o papel do modelo, do domínio e do teto de geração. S207-C1 tem $n = 1$, e S309-C2 mede robustez, não correção.

\paragraph{Planejamento e decomposição}
Avalia-se a geração de plano, pseudocódigo, especificação ou testes intermediários antes do código, num prompt estruturado ou em chamadas encadeadas. Nas sínteses parciais, são cerca de 22 estudos, a maioria com contraste de pacote, além de três estudos de pistas de solução (S892, S1679 e S2145).

% AFIRMACAO: A77
No grupo com raciocínio ou decomposição, 10 contrastes de pacote em 8 estudos registram 8 melhoras e 2 pioras (S1-C1, S6-C1, S77-C1, S216-C1, S216-C2, S295-C2, S491-C1, S654-C1, S654-C2, S845-C1); entre as melhoras estão o Self-planning, de 48,1\% para 60,3\% de Pass@1 no HumanEval com code-davinci-002 contra o prompt direto (S1-C1), e o SCoT, de 53,29\% para 60,64\% contra o CoT (S77-C1).
% AFIRMACAO: A78
As duas pioras ocorreram quando a etapa intermediária veio de um modelo pequeno: a CoT gerada pelo próprio CodeGen 350M baixou o resultado no HumanEval de 14,63\% para 12,8\%, enquanto a de um modelo separado e ajustado o elevou para 20,73\% (S216-C2, S216-C1), e a função auxiliar gerada pelo Incoder 1B mudou o pass@1 de 0,0373 para 0,0364, contra 0,2028 com a função escrita por humanos, condição oráculo (S295-C2, S295-C3).
% AFIRMACAO: A79
Decompor também piorou resultados: com WizardCoder-15B, a geração holística de uma classe Java teve Completion@1 de 98,0\%, contra 74,0\% na síntese independente por método e 70,0\% na incremental, métrica que não mede correção (S365-C1, S365-C3), e o Knowledge Generation Prompting reduziu o pass@1 do GPT-3.5 Turbo no ChemEval de 45,0\% para 39,79\%, em contradição com a narrativa do estudo (S2075-C2).
% AFIRMACAO: A80
Nos dois estudos de pistas de solução com números, os maiores ganhos vieram de informação externa ao modelo avaliado: com DeepSeek-Coder-V2 no MuPA, as pistas autogeradas levaram o pass@1 (escala 0--100) de 11,84 para 13,21, 12,18 e 11,92, e a solução oficial fornecida, para 38,79 (S2145-C1, S2145-C3, S2145-C5, S2145-C6); com Deepseek-coder 33B, os casos de teste aprovados passaram de 269 para 299 com as dicas oficiais e para 801 com dicas reescritas pelo GPT-4o (S1679-C1, S1679-C2).

Os ganhos de pacote não se atribuem à decomposição, e nenhuma comparação mede custo, embora várias técnicas usem mais de uma chamada. Os valores de ``até 25,4\%'' de S1 e ``até 13,79\%'' de S77 são variações relativas máximas, não p.p. As pistas de referência medem o aproveitamento de conhecimento fornecido, não a capacidade de produzi-lo.

\paragraph{Exemplos e sua seleção}
Avaliam-se a presença, a quantidade, a seleção e o formato de demonstrações ou de casos de teste no prompt. Nas sínteses parciais, cerca de 14 estudos do grupo só com instrução, exemplos, contexto ou restrições têm exemplos como técnica focal. Pacotes com seleção treinada ou exploração simbólica (S170-C1, S195-C1, S500-C1) incluem componentes fora do escopo de prompting.

% AFIRMACAO: A81
Quando apenas se acrescentam exemplos, sem recuperação treinada, as diferenças registradas são pequenas: Pass@1 de 0,229 para 0,237 (proporção) no VHDL-Eval (S361-C1), de 59,0\% para 62,8\% no VerilogEval com GPT-4o (S432-C1) e, em R com DeepSeek Coder 1B, de 13,9\% para 14,1\% com \textit{few-shot} e para 13,8\% com exemplos de tradução (S591-C3, S591-C1).
% AFIRMACAO: A82
Os exemplos também pioraram resultados: exemplos de CVT reduziram o pass@1 funcional de 80,81\% para 70,87\% (S1024-C2), amostras de código com erros o derrubaram de 45,0\% para 15,62\% com GPT-3.5 Turbo (S2075-C1), e o \textit{few-shot} elevou a taxa de inacessibilidade do GPT-4o de 0,326 para 0,469 (S2113-C2).

O sinal varia com modelo, temperatura e tarefa (S432, S1707, S2090). Em S902-C2, o conteúdo recuperado recebeu nomes extraídos da solução de referência, condição oráculo, e em S1613-C1 metade dos testes de avaliação foi exposta no prompt.

\paragraph{Contexto recuperado e de repositório}
Avalia-se o acréscimo de documentação de API, código do arquivo ou do repositório, dependências, tipos, dados ou regras de domínio. Nas sínteses parciais, são cerca de 17 estudos do grupo só com instrução, exemplos, contexto ou restrições e 8 do grupo com raciocínio ou decomposição que combinam recuperação e decomposição. Predominam contrastes de pacote; as ablações de S360 e S852 isolam componentes dentro de sistemas completos.

% AFIRMACAO: A83
Com contexto de arquivo ou de repositório, o Pass@1 passou de 17,4\% para 53,04\% no DevEval com gpt-4 (S357-C1) e de 20,29\% para 44,93\% com o RepoScope, que usou em média 3679 \textit{tokens} de prompt contra 97 da base (S852-C1).
% AFIRMACAO: A84
No grupo com raciocínio ou decomposição, 8 das 17 comparações de contexto isolado usam informação oráculo, como o contexto do arquivo somado às APIs invocadas, que elevou o Pass@1 do GPT-4o Mini no RepoExec de 16,62\% para 37,75\% (S213-C2, S213-C3, S213-C5, S295-C3, S670-C2, S670-C5, S670-C6, S670-C7).
% AFIRMACAO: A85
Há inversões e nulos: a recuperação reduziu a acurácia de Claude 3.7 no QHackBench de 31,0\% para 0,0\% (S806-C1), o RAG sobre a documentação do Ansible baixou o BLEU do GPT-4 de 0,1828 para 0,1568 (S1951-C1), a QA API isolada reduziu a acurácia de execução de 59,6\% para 57,4\% (S176-C3) e o InCoder-6B ficou em 0 de Pass@1 em Java com e sem contexto (S137-C1 a S137-C3).

O resultado parece depender do recuperador e da qualidade do que se recupera (S273, S547, S1195). As condições oráculo de S120 e S556 dão limite superior, não condição realista, e há recuperadores treinados em S120, S170 e S852.

\paragraph{Instrução, restrições e persona}
Avaliam-se especificação, pré e pós-condições, restrições de segurança ou acessibilidade, regras de estilo, persona, redação e língua da instrução. Nas sínteses parciais, cerca de 17 estudos do grupo só com instrução, exemplos, contexto ou restrições tratam de instrução e restrições, e 9 de formulação e língua, com contrastes também nos outros grupos.

% AFIRMACAO: A86
Em dois estudos, mudanças só de redação pioraram o resultado, e a persona não ajudou: corrigir problemas de qualidade dos prompts do HumanEval Java reduziu o pass@1 do CodeGen-2.5-7B-multi de 0,219 para 0,15 (S323-C2), cinco de seis frases instrucionais elevaram as soluções incorretas do Codex de 45 para valores entre 59 e 74 em 120, e a sexta as reduziu para 44 (S2162-C1 a S2162-C6), e a persona mudou o Pass@1 do GPT-4o de 47,5\% para 47,4\% (S560-C3).
% AFIRMACAO: A87
Restrições dirigidas a requisitos não funcionais tiveram direção variável: o prompt estruturado de segurança reduziu o RPE do GPT 4o de 0,42 para 0,13 (S1607-C1), mas, com GPT-3.5-turbo, 13 variantes instrucionais sem iteração elevaram a proporção de amostras vulneráveis de 5,74\% para valores entre 5,99\% e 7,82\% (S604-C1 a S604-C10, S604-C13 a S604-C15), as diretrizes OWASP reduziram o Pass@1 de 16,19\% para 9,52\% (S1195-C4) e a instrução WCAG elevou a inacessibilidade de 0,326 para 0,433 (S2113-C1).
% AFIRMACAO: A88
A língua da instrução também alterou o resultado: com StarCoder 3B no X-HumanEval-X, a autotradução da instrução chinesa para o inglês reduziu o pass@1 de 18,29\% para 1,83\%, enquanto a tradução externa o elevou para 20,12\% e 21,95\% (S2141-C1 a S2141-C3), e, com Claude3.5, o prompt em árabe obteve Pass@1 de 0,831 contra 0,938 em inglês (S495-C1).

Scanners aplicados a \textit{benchmarks} não demonstram segurança em produção. O ganho de pré e pós-condições em S1812-C1 é agregado sobre seis modelos e, segundo a ficha, não significativo. Em S739, o objetivo da instrução muda a contagem de vulnerabilidades ao longo de dez iterações, mas sem condição sem iteração.

\paragraph{Otimização automática do prompt}
Avaliam-se a busca iterativa ou evolutiva sobre o prompt, a compressão ou poda e a seleção automática entre estilos. As sínteses parciais registram otimização em S433, S650 e S2053, compressão ou poda em S499, S1644, S1837 e S2048 e seleção em pacotes híbridos em S570, S1853 e S1924.

% AFIRMACAO: A89
Com GPT-3.5, o Prochemy elevou o Pass@1 de 72,6\% para 76,2\% no HumanEval e de 60,5\% para 73,0\% na tradução Java$\rightarrow$Python do CodeNet (S650-C1, S650-C2), e o EPiC elevou o pass@1 do o3-mini no HumanEval+ de 0,88 para 0,945 e 0,941 (proporção) nos dois registros do método (S433-C1, S2053-C1), embora as fichas registrem empate do CoT com e sem Prochemy no LiveCodeBench e o EPiC abaixo do LATS e do LDB no MBPP+ (S650; S2053).
% AFIRMACAO: A90
A compressão da docstring reduziu o Pass@1 do DeepSeekCoder-1.3B no HumanEval de 63,42\% para 54,88\%, 53,05\% e 57,93\%, conforme o método (S499-C1 a S499-C3), enquanto a poda guiada de exemplos, no ClassEval com GPT-3.5-turbo-0613, elevou o Pass@1 de 3,8\% para 9,5\% e de 4\% para 8,6\% contra o \textit{few-shot} sem poda (S1644-C1, S2048-C1).

No ClassEval, os valores absolutos ficam abaixo de 10\%, e em S2048-C1 a intervenção usa mais exemplos que o comparador. A economia de \textit{tokens} de S499 não equivale a redução de custo medida. Os pacotes híbridos incluem seleção de modelo ou de \textit{solver}, e S1853 mede o tempo de busca; nenhum desses ganhos se atribui à seleção de prompt.

\paragraph{Estudos extraídos depois das sínteses parciais}
S1610 e S1617 entraram depois da decisão do autor sobre a data de I4 e não estão nas contagens dos agrupamentos
acima.
% AFIRMACAO: A103
Em S1610, com Qwen 2.5 Coder 7B e 81 prompts, a proporção de arquivos gerados com ao menos um alerta do Bandit caiu
de 63,0\% no zero-shot para 54,3\% com \textit{few-shot} e para 55,6\% com CoT. Nas duas condições, o contraste também
acrescenta uma diretiva de segurança, e por isso é de pacote (S1610-C1; S1610-C2).
% AFIRMACAO: A104
Em S1617, com GPT-5, Gemini 2.5 Pro e DeepSeek V3.2 no HumanEval-XL, o Pass@1 agregado foi de 0,858 com prompts em inglês e de 0,844 com prompts em espanhol (S1617-C1).

\paragraph{Lacunas}
As sínteses parciais deixam lacunas que limitam os achados de geração:
\begin{itemize}
\item \textbf{Isolamento com orçamento igual.} Só S479 e S645 igualam o orçamento do laço de \textit{feedback}; nenhuma comparação contrasta multiagente e agente único com as mesmas chamadas, nem separa raciocínio de decomposição no mesmo pacote.
\item \textbf{Custo e uso real.} Nenhuma comparação mede custo ou tempo de forma pareada, e os \textit{benchmarks} não sustentam conclusões sobre produtividade, segurança ou redução de custo.
\item \textbf{Informação oráculo.} Condições com informação da solução de referência (S120, S295-C3, S399-C3, S556, S670, S892, S902-C2, S1613-C1, S1708, S2140-C2, S2145-C6) indicam limites superiores.
\item \textbf{Modelos e domínios.} Modelos pequenos (por exemplo, S114, S479, S759 e S1977) e linguagens de hardware ou de domínio aparecem em poucos estudos, com conjuntos próprios.
\item \textbf{Comparações de uma célula.} Cada comparação cobre um modelo e um \textit{benchmark}; as inversões entre modelos ficam só como heterogeneidade.
\item \textbf{Qualidade e conferência.} A contaminação foi julgada ``atende'' em 9 dos 222 estudos das sínteses parciais (S247, S323, S397, S398, S560, S782, S902, S2054 e S2123), e nenhuma ficha passou por conferência humana. S433 e S2053 descrevem o EPiC com o mesmo modelo e \textit{benchmark}; esta revisão não verificou se são versões do mesmo trabalho, o que afetaria as contagens.
\end{itemize}
